
\source{Лекции, с. 1–3, 10, 13–15, 19–20; обозначения дополнены для билетов}
\begin{definition}[Граф и подграф]
Неориентированный граф \(G=(V,E)\) состоит из конечного непустого множества
вершин \(V\) и множества неупорядоченных пар различных вершин \(E\).
Таким образом, по умолчанию петель и кратных рёбер нет.
Орграф задаётся множеством дуг \(E\subseteq V^2\); дуга \((u,v)\)
направлена из \(u\) в \(v\). Забывая направления дуг, удаляя петли и кратности,
получаем соответствующий неориентированный граф.
Подграф \(H=(W,F)\) удовлетворяет \(W\subseteq V\), \(F\subseteq E\),
причём концы всех рёбер \(F\) лежат в \(W\).
Подграф индуцирован множеством \(W\), если содержит все рёбра исходного
графа с обоими концами в \(W\); он обозначается \(G[W]\).
Остовный подграф содержит все вершины: \(W=V\).
\end{definition}
Для орграфа направление каждой дуги учитывается. Иные соглашения оговорены в билетах.
Пишем \(G=(V,E)\), \(n=|V|\), \(m=|E|\), \(c(G)\) — число компонент связности,
\(N(v)\) — множество соседей, \(d(v)=\deg(v)=|N(v)|\),
\(\delta(G)=\min d(v)\), \(\Delta(G)=\max d(v)\).
\(G-v\) и \(G-e\) — удаление вершины со всеми инцидентными рёбрами и удаление ребра;
\(G[S]\) — индуцированный подграф на \(S\).
\(K_n\), \(K_{p,q}\), \(P_n\), \(C_n\) — полный граф, полный двудольный граф,
простой путь и простой цикл соответственно; в \(P_n,C_n\) индекс означает число вершин.

\textbf{Путь} (в терминологии лекций; также \textbf{маршрут}) — последовательность \(v_0,e_1,v_1,\ldots,e_k,v_k\),
где каждое ребро соединяет соседние вершины последовательности.
Его длина равна \(k\); разрешены повторения вершин и рёбер.
В орграфе дуги проходятся по направлению.
\textbf{Цепь} — маршрут без повторений рёбер; \textbf{простой путь} — маршрут
без повторений вершин. Путь с одинаковыми началом и концом называется замкнутым.
В лекционном определении неориентированного \textbf{цикла} дополнительно
требуются положительная длина и отсутствие немедленного возвращения
по тому же ребру, в том числе на стыке конца и начала записи.
Поэтому путь \(u,v,u\) по одному ребру замкнут, но циклом не считается.
Замкнутая цепь дополнительно не повторяет рёбер;
\textbf{простой цикл} — замкнутый маршрут положительной длины без других
повторений вершин (в простом неориентированном графе длины не менее трёх).

Записи неориентированного цикла, отличающиеся циклическим сдвигом
или обращением обхода, описывают один цикл. Для ориентированного цикла
обращение допускается только при наличии нужных дуг; при подсчёте замкнутых
путей с отмеченным началом сдвиги считаются разными записями.
Цикл может проходить один и тот же цикл несколько раз; простой — нет.
Если \(P=R^r\), \(Q=R^s\) для положительных \(r,s\), где степень означает
склейку повторных обходов, то в лекциях такие циклы также отождествляются
отдельным отношением \(P\equiv Q\) (после сдвига/обращения записи).
Это отождествление не применяется при подсчёте путей заданной длины.

\textbf{Лемма (лекции, с. 15).} Наличие цикла в неориентированном графе
влечёт наличие простого цикла. Выбираем в записи цикла кратчайший
непустой замкнутый участок между повторениями вершины. Внутренних
повторений у него нет по минимальности; длина 1 исключена отсутствием
петель, длина 2 — запретом немедленного возврата по ребру.
Следовательно, участок — простой цикл.

Из любого пути между различными вершинами можно получить простой путь:
при повторении вершины вырезать участок между двумя её вхождениями.
Длина строго уменьшается, поэтому процесс заканчивается.
Граф \textbf{связен}, если любые две его вершины соединены путём.
Связность является отношением эквивалентности: нулевой путь даёт
рефлексивность, обращение пути — симметричность, склейка — транзитивность.
Её классы — компоненты связности.
Одновершинный граф связен. При удалениях допускается пустой граф:
у него ноль компонент. Это вспомогательное расширение лекционного соглашения.

\textbf{Лемма о рукопожатиях.} \(\sum_{v\in V}d(v)=2m\).
Действительно, каждое ребро учитывается ровно по одному разу у каждого конца.
В частности, число вершин нечётной степени чётно.
Для орграфа принимаем знаки со с. 14 лекций:
\(d^+(v)=|\{u:uv\in E\}|\) — число \emph{входящих},
\(d^-(v)=|\{u:vu\in E\}|\) — число \emph{исходящих} дуг.
В других источниках знаки часто противоположны; здесь везде используется
указанное лекционное соглашение.
\(\sum d^+(v)=\sum d^-(v)=|E|\), поскольку каждая дуга
имеет ровно одно начало и один конец.

\textbf{Линейная алгебра.} Используются определитель, его разложение по строке,
мультилинейность и правило умножения матриц. Необходимая для билета 7
формула Бине–Коши обосновывается непосредственно там.
Для плоских графов используется теорема Жордана: простая замкнутая кривая
разделяет плоскость на внутреннюю и внешнюю области и служит их общей границей.
Это топологический исходный факт, а не дополнительная графовая лемма.

\textbf{Обозначения матриц и раскрасок.}
Матрица смежности — \(A(G)\), матрица Кирхгофа — \(K(G)\).
Хроматический многочлен обозначается \(\chi_G(t)\), как в лекциях;
его аргумент может обозначаться также \(k\). Хроматическое число —
\(\chi(G)\), без нижнего индекса: это число, а не многочлен.
